\documentclass[10pt,a4paper]{amsart}
\usepackage[margin=19mm]{geometry}
\usepackage{amsmath,amssymb,mathtools}
\usepackage[scr=rsfs,cal=euler]{mathalfa}
\usepackage{xfrac}
\usepackage[unicode,hidelinks,bookmarks=false]{hyperref}
\newcommand{\ie}{i.e.\ }
\newcommand{\cf}{cf.\ }
\newcommand{\resp}{respectively\ }
\newcommand{\Subsection}{\subsection}
\providecommand{\nobreakdash}{\nobreak\hbox{-}}
\setlength{\emergencystretch}{2em}
\multlinegap0pt
\usepackage{latexsym,mathrsfs,verbatim}
\newtheorem{theorem}{Theorem}[section]
\newtheorem{corollary}[theorem]{Corollary}
\newtheorem{lemma}[theorem]{Lemma}
\newtheorem{proposition}[theorem]{Proposition}
\newtheorem{remark}[theorem]{\bf Remark}
\newcommand{\cA}{\mathcal A}
\newcommand{\cP}{\mathcal P}
\newcommand{\cC}{\mathcal C}
\newcommand{\LL}{\pounds}
\newcommand{\leg}[2]{\left(\frac{#1}{#2}\right)}
\DeclareMathOperator{\logA}{\log_{\mathcal{A}}}
\DeclareMathOperator{\logAm}{\log_{\mathcal{A}(m)}}
\DeclareMathOperator{\Pm}{{\mathcal P}(\textrm{m})}
\DeclareMathOperator{\ord}{ord}
\DeclareMathOperator{\rad}{rad}
\DeclareMathOperator{\Rec}{Rec}
\DeclareMathOperator{\im}{im}
\DeclareMathOperator{\Gal}{Gal}
\DeclareMathOperator{\Frob}{Frob}
\DeclareMathOperator{\Tr}{Tr}
\begin{document}
\title[Irrationality of finite logarithms in a congruence-class ad`]{Irrationality of finite logarithms in a congruence-class ad\`ele ring}
\author{Daniel Evans}
\date{}
\maketitle
\noindent\emph{Source and licence.} Irrationality of finite logarithms in a congruence-class ad\`ele ring. Version arXiv:2607.27774v2 (CC BY 4.0). This material is distributed under the Creative Commons Attribution 4.0 International licence, http://creativecommons.org/licenses/by/4.0/, as recorded in the arXiv OAI licence metadata for this exact version and shown on the abstract page https://arxiv.org/abs/2607.27774v2. Full rights notice: licenses/arxiv-2607.27774-CC-BY-4.0.txt.\par\smallskip
\noindent\textit{Editorial scope.} Counted unit: the original \texttt{theorem} environment \texttt{logAm} at source lines 89--97 together with its complete proof at lines 244--311; the delivered document keeps source lines 53--433 so that every referenced label and every citation resolves inside it. Native editable source is the paper's own concatenated TeX; the AMS wrapper is editorial formatting only. No publisher class, upstream script or unrelated artwork is redistributed.
\section{Introduction}
In recent years, there has been growing interest in arithmetic in the so-called {\it poor man's ad\`{e}le ring}
\[\cA:=\left.\left(\;\prod_{p\in\cP}\;\mathbb{Z}/p\mathbb{Z}\right) \middle/ \left(\;\bigoplus_{p\in\cP}\;\mathbb{Z}/p\mathbb{Z}\right)\right.. \]
The direct product and sum are taken over the set of primes $\cP$, and addition and multiplication are taken component-wise.
The field of rational numbers $\mathbb{Q}$ embeds naturally along the diagonal, and in this manner we can view $\mathcal{A}$ as a $\mathbb{Q}$-algebra. 
Explicitly, elements ${\bf t}\in\cA$ can be encoded as vectors
\[{\bf t}=(t_2, t_3, t_5,...) = (t_p)_{p\in\mathcal{P}}\quad\text{where $t_p\in\mathbb{Q}$}\]
and a finite number of reductions $t_p\bmod p$ may remain undefined when $t_p$ is not $p$-integral. 
One then identifies ${\bf t}={\bf s}$ when $t_p\equiv s_p\bmod p$ for all sufficiently large primes $p\in\cP$. Going further, Rosen (see \cite{Ro20} and \cite{RTTY24}) has developed a notion of algebraic numbers in $\mathcal{A}$, which are essentially ${\bf t}=(t_p)_{p\in\mathcal{P}}$ arising from sequences of rational numbers $t_n$ satisfying a linear recurrence with constant rational coefficients. 
Alongside many aspects concerning finite multiple zeta values, in \cite{KZ} Kaneko and Zagier discuss several other arithmetically interesting special elements in $\cA$, including Fermat quotients modulo $p$ as the prime $p$ varies. 
For prime $p$, write $\mathbb{Z}_{(p)}$ for the localisation of $\mathbb{Z}$ at the prime ideal $(p)$. 
Given $\alpha\in\mathbb{Z}_{(p)}^\times$, the {\it Fermat quotient} is
\[q_p(\alpha)=\frac{\alpha^{p-1}-1}{p}\in\mathbb{Z}_{(p)}\]
and we have a corresponding function
\begin{align*}
\logA:\mathbb{Q}^\times & \longrightarrow\cA \\
\alpha& \longmapsto \left(q_p(\alpha)\bmod p\right)_{p\in\cP}.
\end{align*}
It is easy to check $\log_\cA(\alpha\beta)=\log_\cA(\alpha)+\log_\cA(\beta)$, so this provides a natural analogue of the logarithm function with values in $\cA$. On investigating irrationality and transcendence in $\cA$, Luca and Zudilin \cite{LZ25} suggested an interesting step would be to prove $\log_\cA(2)$ is irrational. This was answered by Matsusaka and Seki \cite{MS26}, who went further, showing that $\log_\cA(\alpha)$ cannot take non-zero rational values in $\cA$ for $\alpha\in\mathbb{Q}^\times\setminus\{\pm 1\}$. Furthermore, an earlier result of Silverman \cite{Si88} on non-Wieferich primes shows that $\log_\cA(\alpha)$ cannot be zero, provided the $abc$-conjecture holds.
\begin{theorem}[Matsusaka--Seki, Silverman]\label{MSS}
For any $\alpha\in\mathbb{Q}^\times\setminus\{\pm 1\}$, we have 
\begin{itemize}
\item[(i)]
$\log_\cA(\alpha)\in\cA\setminus\mathbb{Q}^{\times}$, 
\item[(ii)]
$\log_\cA(\alpha)\neq 0\in\cA$ assuming the $abc$-conjecture.
\end{itemize}
\end{theorem}
In \cite{GRM13}, Graves and Ram Murty extended Silverman's method to prove Theorem \ref{MSS}(ii) for integer $\alpha>1$ when primes are restricted to congruence classes 
\[\cP(m)=\{p\in\cP\;|\; p\equiv 1\bmod m\}\]
for fixed integer $m\geq 1$. In this note, we will extend this to rational $\alpha$ and give a similar extension to Theorem \ref{MSS}(i) for these congruence classes. 
To state this, define a corresponding restricted poor man's ad\`{e}le ring
\[\cA(m):=\left.\left(\;\prod_{p\in\cP(m)}\mathbb{Z}/p\mathbb{Z}\right) \middle/ \left(\;\bigoplus_{p\in\cP(m)}\mathbb{Z}/p\mathbb{Z}\right)\right..\]
The set $\cP(m)$ is infinite with density $1/\varphi(m)$ in $\cP$ by Chebotarev's density theorem, where $\varphi(m)$ denotes Euler's totient function. 
Again, the field of rational numbers $\mathbb{Q}$ embeds in $\cA(m)$ naturally, and we consider the corresponding logarithm 
\[\log_{\cA(m)}(\alpha)=(q_p(\alpha)\bmod p)_{p\in\cP(m)}.\]

\begin{theorem}\label{logAm}
For any $\alpha\in\mathbb{Q}^\times\setminus\{\pm 1\}$ and $m\geq 1$, we have
\begin{itemize}
\item[(i)]
$\log_{\cA(m)}(\alpha)\in\cA(m)\setminus\mathbb{Q}^{\times}$, 
\item[(ii)]
$\log_{\cA(m)}(\alpha)\neq 0\in\cA(m)$ assuming the $abc$-conjecture.
\end{itemize}
\end{theorem}

Theorem \ref{MSS} can be rephrased in the following way: for fixed $\alpha\in\mathbb{Q}^\times\setminus\{\pm 1\}$ and $c\in\mathbb{Q}$, there are infinitely many primes $p$ such that 
\[q_p(\alpha)=\frac{\alpha^{p-1}-1}{p}\not\equiv c\bmod p\]
where the case $c=0$ is dependent on the $abc$-conjecture. Theorem \ref{logAm} says there are infinitely many primes satisfying this condition in arithmetic progressions of the form $p\equiv 1\bmod m$. 
The methods in both \cite{MS26} and \cite{Si88} revolve around finding suitable congruences and bounds for values of cyclotomic polynomials $\Phi_\ell(\alpha)$ with sufficiently large primes $\ell$, assuming $\log_\cA(\alpha)$ is rational. We modify this by considering $\Phi_{m\ell}(\alpha)$, together with a suitable logarithmic derivative correction.
\smallskip

The machinery used to prove Theorem \ref{logAm} can be adapted to give further results about $\log_\cA(\alpha)$ in $\cA$. 
Rather than excluding a single non-zero rational value $q_p(\alpha)=c$, a signed version of the cyclotomic argument allows us to exclude the possibility that $q_p(\alpha)=\pm c$ simultaneously for $c\in\mathbb{Q}^\times$. 
In particular, we obtain a second strengthening of Theorem \ref{MSS}.
\begin{theorem}\label{logsquare}
For any $\alpha\in\mathbb{Q}^\times\setminus\{\pm 1\}$, we have 
\begin{itemize}
\item[(i)]
$\log_\cA(\alpha)^2\in\cA\setminus\mathbb{Q}^{\times}$, 
\item[(ii)]
$\log_\cA(\alpha)^2\neq 0\in\cA$ assuming the $abc$-conjecture.
\end{itemize}
\end{theorem}

The restricted congruence class result also has an application concerning $\log_\cA(\alpha)$ in Rosen's theory of finite algebraic numbers. Following a suggestion by Luca and Zudilin in \cite{LZ25}, we define the recurrence degree $\deg_\cA({\bf t})$ of ${\bf t}\in\cA$ to be the least order of a rational constant-coefficient recurrence whose $p$-th terms represent ${\bf t}$. We obtain the following consequence of Theorem \ref{logAm}.
\begin{theorem}\label{degtwo}
Let $\alpha\in\mathbb{Q}^\times\setminus\{\pm 1\}$. If $\deg_\cA(\log_\cA(\alpha))\leq 2$, 
then there exists $r\in\mathbb{Q}$ and a square-free integer $D\neq 1$ such that
\[\log_\cA(\alpha)=r\left(1-\leg{D}{p}\right)_p.\]
Furthermore, if the $abc$-conjecture holds, then $\deg_\cA(\log_\cA(\alpha))>2$. 
\end{theorem}
Section 2 contains necessary cyclotomic constructions and we give the proofs of Theorem \ref{logAm} and Theorem \ref{logsquare} in Section 3. In Section 4, we introduce the recurrence degree, classify elements of degree at most 2, and use this to prove Theorem \ref{degtwo}. 


 
\section{Cyclotomic preliminaries}
For integer $n\geq 1$, let $\zeta_n$ be a primitive $n$th root of unity and $\Phi_n(T)$ be the $n$th cyclotomic polynomial with degree $\varphi(n)$
\begin{equation*}
\Phi_n(T)=\prod_{\substack{1\leq j<n \\ \gcd(j,n)=1}} \left(T-\zeta_n^j\right)\in\mathbb{Z}[T].
\end{equation*}
We will also make use of the homogenised form
\begin{equation}\label{homPhi}
\Phi_{n}(X,Y)=Y^{\varphi(n)}\Phi_n(X/Y)=
\prod_{\substack{1\leq j<n \\ \gcd(j,n)=1}} \left(X-\zeta_n^jY\right)\in\mathbb{Z}[X,Y]
\end{equation}
as well as its partial derivatives $\Phi_{n,X}(X,Y)$ and $\Phi_{n,Y}(X,Y)$. By Euler's homogeneous form identity, these satisfy
\begin{equation}\label{euler}
X\Phi_{n,X}(X,Y)+Y\Phi_{n,Y}(X,Y)=\varphi(n)\Phi_n(X,Y).
\end{equation}
Given $n\geq 1$ and a prime $p\nmid n$, the well-known identity $\Phi_n(T^p)=\Phi_n(T)\Phi_{np}(T)$ has homogenised version
\begin{equation}\label{Phinp}
\Phi_n(X^p,Y^p)=\Phi_n(X,Y)\Phi_{np}(X,Y). \\[10pt]
\end{equation}
To investigate $q_p(\alpha)$ for $\alpha=u/v\in\mathbb{Q}^\times\setminus\{\pm 1\}$, since 
\[q_p(\alpha)\equiv q_p(-\alpha)\equiv -q_p(\alpha^{-1})\bmod p\]
for $p>2$, we need only consider $u>v\geq 1$ with $\gcd(u,v)=1$, and will assume this from here onwards. 
We first collect some standard facts regarding prime factors of $\Phi_n(u,v)$. 
\begin{lemma}\label{ordlem}
Given prime $p$ dividing $\Phi_n(u,v)$, we have the following.
\begin{itemize}
\item[(i)]
$p\nmid uv$,
\item[(ii)]
if $p\nmid n$, then $\ord_p(\alpha)=n$ and $p\equiv 1\mod n$,
\item[(iii)]
if $p\mid n$, then $\ord_p(\alpha)=n/p^k$ for some $k\geq 1$ and $p$ is the largest prime factor of $n$.
\end{itemize}
\end{lemma}
\begin{proof}
Since $\Phi_n(u,v)\mid \left(u^n-v^n\right)$, we have $u^n\equiv v^n\bmod p$, so if $p$ divided one of $u$ or $v$, it would divide both. This contradicts $\gcd(u,v)=1$ and part (i) follows.
\medskip

The reduction $\alpha\bmod p$ is a well-defined element in $\mathbb F_p^\times$ and $p$ divides $\Phi_n(u,v)=v^{\varphi(n)}\Phi_{n}(\alpha)$,  thus $\Phi_{n}(\alpha)\equiv 0\bmod p$. 
Writing $n=n_0p^k$ where $p\nmid n_0$, the following identity in $\mathbb{F}_p[T]$
\[\Phi_n(T)\equiv\Phi_{n_0}(T)^{\varphi(p^k)}\bmod p\]
can be proved easily using \eqref{Phinp} or via M\"{o}bius inversion. Thus $\Phi_{n_0}(\alpha)\equiv 0\bmod p$ and we have $\alpha^{n_0}\equiv 1\bmod p$. 
Suppose that $\alpha^d-1\equiv 0\bmod p$ for some $d\mid n_0$. Then $\alpha$ is a root modulo $p$ of 
\[T^{d}-1=\prod_{e\mid d}\Phi_e(T)\]
and hence $\Phi_e(\alpha)\equiv 0\bmod p$ for some $e\mid d$ and $d\mid n_0$. 
However, the polynomial $T^{n_0}-1$ is separable over $\mathbb{F}_p$ and the factors in the decomposition
\[T^{n_0}-1=\prod_{d\mid n_0}\Phi_d(T)\]
are pairwise coprime over $\mathbb{F}_p$. We deduce that $e=d=n_0$ and so $\ord_p(\alpha)=n_0=n/p^k$. 
\medskip

Clearly this order $n_0$ divides $\left|\mathbb{F}_p^\times\right|=p-1$ and when $k=0$, we obtain $p\equiv 1\bmod n$. 
When $k\geq 1$, any prime $q\neq p$ dividing $n$ divides $n_0$ and so $q\mid p-1$. Hence $p$ is the largest prime factor of $n$.
\end{proof}

We next relate Fermat quotients with values of cyclotomic polynomials. 
\begin{lemma}\label{conglem}
Suppose that prime $p\mid\Phi_n(u,v)$ and $p\nmid n$. Then
\[\dfrac{\Phi_{n}(u,v)}{p}\equiv q_p(\alpha)v\Phi_{n,Y}(u,v)\bmod p\] 
and moreover, $\Phi_{n,Y}(u,v)\not\equiv 0\bmod p$.
\end{lemma}
\begin{proof} 
Using \eqref{euler} and the assumption $p\mid\Phi_n(u,v)$, we have
\begin{equation}\label{euler2}
u\Phi_{n,X}(u,v)\equiv -v\Phi_{n,Y}(u,v)\bmod p.
\end{equation}
Now, $u^p-u\equiv v^p-v\equiv 0\bmod p$, so by taking a Taylor expansion we find
\begin{equation}\begin{aligned}[b]
\Phi_n(u^p,v^p) 
&\equiv \Phi_n(u,v)+(u^p-u)\Phi_{n,X}(u,v)+(v^p-v)\Phi_{n,Y}(u,v)\bmod p^2  \\
&\equiv \Phi_n(u,v)+p\left[q_p(u)u\Phi_{n,X}(u,v)+q_p(v)v\Phi_{n,Y}(u,v)\right]\bmod p^2 \\
&\equiv\Phi_n(u,v)-p\left(q_p(u)-q_p(v)\right)v\Phi_{n,Y}(u,v) \bmod p^2 \\
&\equiv\Phi_n(u,v)-pq_p(\alpha)v\Phi_{n,Y}(u,v)\bmod p^2. \label{cong}
\end{aligned}\end{equation}
Note that since $p\nmid n$, we have
\[\Phi_n(X,Y)\Phi_{np}(X,Y)=\Phi_n(X^p,Y^p)\equiv\Phi_n(X,Y)^p\bmod p.\]
and so $\Phi_{np}(X,Y)\equiv\Phi_n(X,Y)^{p-1}\bmod p$. 
In particular, $\Phi_{np}(u,v)\equiv\Phi_{n}(u,v)^{p-1}\equiv 0\bmod p$ and hence 
\[\Phi_n(u^p,v^p)=\Phi_n(u,v)\Phi_{np}(u,v)\equiv 0 \bmod p^2.\]
The congruence follows on dividing \eqref{cong} by $p$.
\medskip

By the previous lemma, $\alpha$ is a simple root of $\Phi_n(T)\bmod p$ so we have $\Phi_n'(\alpha)\not\equiv 0\bmod p$. Hence
\[\Phi_{n,X}(u,v)=v^{\varphi(n)-1}\Phi_n'(\alpha)\not\equiv 0\bmod p\] 
and $\Phi_{n,Y}(u,v)\not\equiv 0\bmod p$ follows from \eqref{euler2}.
\end{proof}
\medskip

Finally, we need a lemma controlling primes $p\mid\Phi_n(u,v)$ as $n$ increases. We first record the trivial estimates
\begin{equation}\label{trivest}
(u-v)^{\varphi(n)}\leq\Phi_n(u,v)\leq (u+v)^{\varphi(n)}
\end{equation}
which follow from the product expansion in \eqref{homPhi}. Note we are assuming $u>v\geq 1$ throughout so $\Phi_n(u,v)$ is positive.
\begin{lemma}\label{limlem}
For $n>1$, let $r_n=\#\{p\text{ prime} : p\mid\Phi_n(u,v), p\nmid n\}$. Then 
\[r_n<\frac{\varphi(n)\log(u+v)}{\log n}
\qquad\text{and}\qquad
\lim_{n\to\infty}\sum_{\substack{p\mid\Phi_n(u,v) \\ p\nmid n}}\frac{1}{p}=0.\]
\end{lemma}
\begin{proof}
When $r_n=0$, the inequality holds trivially so assume $r_n>0$. For each prime $p$ dividing $\Phi_n(u,v)$ with $p\nmid n$, 
we have $p\equiv 1\mod n$ by Lemma \ref{ordlem}, so $p>n$. By \eqref{trivest},
\begin{align*}
n^{r_n}<\prod_{\substack{p\mid\Phi_n(u,v) \\ p\nmid n}}\!\! p\;\ \leq \Phi_n(u,v)
\leq (u+v)^{\varphi(n)}
\end{align*}
and the inequality for $r_n$ follows. Furthermore, $\varphi(n)<n$ so
\[\sum_{\substack{p\mid\Phi_n(u,v) \\ p\nmid n}}\frac{1}{p}
\leq\frac{r_n}{n}
<\frac{\log(u+v)}{\log n}\longrightarrow 0\qquad\text{as $n\to\infty$.}\]
\vspace{-1cm}

\end{proof}


\section{Irrationality of finite logarithms}
\subsection{Non-zero rational values in $\cA(m)$} $\,$\\[3pt]
To prove Theorem \ref{logAm}(i), we adapt the method used in \cite{MS26} to fit the template of many classical irrationality proofs by contradiction: assuming rationality, define a sequence of non-zero integers that tends to zero.

\begin{proof}[Proof of Theorem \ref{logAm}(i)] 
Fix $m\geq 1$ and suppose for contradiction there exist $a, b\in\mathbb{Z}$ with $a\neq 0$, $b\geq 1$, $\gcd(a,b)=1$ and a bound $N>\max\{|a|,b\}$ such that for all primes $p>N$ with $p\equiv 1\bmod m$,  we have
\[q_p(\alpha)\equiv \frac{a}{b}\bmod p.\] 
\vspace{-0.1cm}

\noindent
Let $w=\Phi_m(u,v)$ and consider $W_\ell=\Phi_{m\ell}(u,v)$ for primes $\ell>\max\{N, m, (u+v)^{\varphi(m)}\}$. 
Since $\ell\nmid m$, by \eqref{Phinp} we have  
\[W_\ell=\Phi_{m\ell}(u,v)=\frac{\Phi_m(u^\ell,v^\ell)}{w}.\]
We also note $\ell>w$ since $(u+v)^{\varphi(m)}\geq w$ by \eqref{trivest}.
\vspace{0.2cm}

Given a prime $p\mid W_\ell$, we automatically have $p\nmid m\ell$. Indeed, if $p\mid m\ell$, then by Lemma \ref{ordlem}(iii) the largest prime factor of $m\ell$ would be $\ell=p$. However, 
\[wW_\ell=\Phi_m(u^\ell,v^\ell)\equiv\Phi_m(u,v)^\ell\equiv w^\ell\mod\ell\]
and so $W_\ell\equiv w^{\ell-1}\equiv 1\mod\ell$, contradicting $\ell=p\mid W_\ell$. 
Consequently, $p\equiv 1\mod m\ell$ by Lemma \ref{ordlem}(ii) so $p\equiv 1\mod m$ and $p>\ell>N$.
Hence $q_p(\alpha)\equiv ab^{-1}\not\equiv 0\bmod p$ and Lemma \ref{conglem} therefore gives
\begin{equation*}\label{Wlcong}
\frac{W_\ell}{p}\equiv q_p(\alpha)v\Phi_{m\ell,Y}(u,v)\not\equiv 0\bmod p.
\end{equation*}
This means $p^2\nmid W_\ell$ for every prime $p\mid W_\ell$ and so $W_\ell$ is square-free.
\bigskip

Differentiate $\Phi_m(X^\ell,Y^\ell)=\Phi_m(X,Y)\Phi_{m\ell}(X,Y)$ with respect to $Y$ to get
\[\ell Y^{\ell-1}\Phi_{m,Y}(X^\ell,Y^\ell)=\Phi_m(X,Y)\Phi_{m\ell,Y}(X,Y)+\Phi_{m\ell}(X,Y)\Phi_{m,Y}(X,Y).\]
Using this, define the following integer
\[R_\ell=\ell v^\ell\Phi_{m,Y}(u^\ell,v^\ell)=wv\Phi_{m\ell,Y}(u,v)+W_\ell v\Phi_{m,Y}(u,v)\]
and for each $p\mid W_\ell$, we now have a congruence
\begin{equation*}\label{wWlcong}
w\frac{W_\ell}{p}\equiv q_p(\alpha)R_\ell\bmod p.
\end{equation*}
Given that $q_p(\alpha)\equiv ab^{-1}\bmod p$, for each prime $\ell$ as above, set
\begin{equation}\label{Tell}
T_\ell=bwS_\ell-aR_\ell\in\mathbb{Z}\qquad\text{where}\qquad S_\ell=\sum_{p\mid W_\ell}\frac{W_\ell}{p}\in\mathbb{Z}.
\end{equation}
Notice that for each specific prime $p\mid W_\ell$, the summands $W_\ell/p'$ for $p'\neq p$ are divisible by $p$ so
$S_\ell\equiv W_\ell/p\mod p$. In particular, 
\[T_\ell\equiv bw\frac{W_\ell}{p}-aR_\ell\equiv bq_p(\alpha)R_\ell-aR_\ell\equiv 0\bmod p.\]
Thus primes dividing $W_\ell$ also divide $T_\ell$, and since $W_\ell$ is square-free, we obtain $T_\ell/W_\ell\in\mathbb{Z}$.
\bigskip

We next show the sequence $T_\ell/W_\ell$ tends to zero. By Lemma \ref{limlem}, the sum $\sum_{p\mid W_\ell}1/p\to 0$ as $\ell\to\infty$. Also,
\[\frac{R_\ell}{W_\ell}=\ell wv^\ell\frac{\Phi_{m,Y}(u^\ell,v^\ell)}{\Phi_m(u^\ell,v^\ell)}
=\ell w\alpha^{-\ell}\frac{\Phi_{m,Y}(1,\alpha^{-\ell})}{\Phi_m(1,\alpha^{-\ell})}
\]
where we have used that $\Phi_m(X,Y)$ is homogeneous of degree $\varphi(m)$ and $\Phi_{m,Y}(X,Y)$ is homogeneous of degree $\varphi(m)-1$. The logarithmic derivative factor is bounded. Indeed, taking the logarithmic $Y$-derivative of \eqref{homPhi} with $n=m$ and setting $X=1$ and $Y=\alpha^{-\ell}$, we have 
\[\left|\frac{\Phi_{m,Y}(1,\alpha^{-\ell})}{\Phi_m(1,\alpha^{-\ell})}\right|\leq
\sum_{\substack{1\leq j<m \\ \gcd(j,m)=1}}\left|\frac{-\zeta_m^j}{1-\zeta_m^j\alpha^{-\ell}}\right|
\leq\frac{\varphi(m)}{1-\alpha^{-\ell}}.\]
Hence 
\[\left|\frac{R_\ell}{W_\ell}\right|\leq\ell w\varphi(m)\frac{\alpha^{-\ell}}{1-\alpha^{-\ell}}=\frac{lw\varphi(m)}{\alpha^\ell-1}\]
and this tends to zero since $\alpha=u/v>1$. Overall, 
\[\left|\frac{T_\ell}{W_\ell}\right|\leq bw\sum_{p\mid W_\ell}\frac{1}{p}+|a|\left|\frac{R_\ell}{W_\ell}\right|\longrightarrow 0
\]
as primes $\ell\to\infty$.
\medskip

Finally, we show $T_\ell\not\equiv 0\mod \ell$. 
As before, we have $W_\ell\equiv 1\bmod \ell$ and each prime divisor $p$ of $W_\ell$ satisfies $p\equiv 1\mod \ell$ and hence $W_\ell/p\equiv 1\mod\ell$. By definition, $\ell\mid R_\ell$ so
\[T_\ell\equiv bwr_{m\ell}\bmod \ell\]
using the notation from Lemma \ref{limlem}. We have chosen $\ell$ so that $1\leq b<\ell$ and $1\leq w<\ell$ so $bw\not\equiv 0\bmod\ell$. We also chose $\ell>(u+v)^{\varphi(m)}$ and hence by Lemma \ref{limlem}, 
\[r_{m\ell}<\frac{(\ell-1)\varphi(m)\log(u+v)}{\log(m\ell)}<\ell-1.\]
Furthermore, by \eqref{trivest}, 
\[W_\ell=\frac{\Phi_m(u^\ell,v^\ell)}{\Phi_m(u,v)}\geq\left(\frac{u^\ell-v^\ell}{u+v}\right)^{\varphi(m)}>1
\]
since $u^\ell-v^\ell>u+v$ for every $\ell\geq 3$. Thus $W_\ell$ has a prime divisor and $1\leq r_{m\ell}<\ell$. 
We conclude that $T_\ell\not\equiv 0\bmod \ell$. In particular, $T_\ell/W_\ell$ is a non-zero integer and we obtain the required contradiction.
\end{proof}


\subsection{Zero values in $\cA(m)$} $\,$\\[3pt]
Graves and Ram Murty (Theorem 3.1 in \cite{GRM13}) demonstrate that there are infinitely many primes $p\equiv 1\bmod m$ such that $q_p(\alpha)\not\equiv 0\bmod p$ for integer $\alpha>1$, assuming the $abc$-conjecture. In fact, they show that the number of such primes not exceeding $x$ is $\gg \frac{\log x}{\log\log x}$ and this bound was later improved to $\gg\log x$ by Ding \cite{D19}. For completeness, we will briefly show there are infinitely many such primes for arbitrary rational $\alpha\in\mathbb{Q}^\times\setminus\{\pm 1\}$. 

The input from the $abc$-conjecture appears via the following result used in Silverman's proof of Lemma 7 in \cite{Si88}.
\begin{lemma}\label{Silabc}
Let $\alpha=u/v$ where $u>v\geq 1$ with $\gcd(u,v)=1$. For each $n\geq 1$, write $u^n-v^n=A_nB_n$ where $B_n$ is the powerful part of $u^n-v^n$, that is
\[A_n=\prod_{p||u^n-v^n} p\qquad\text{and}\qquad 
  B_n=\prod_{\substack{p^e||u^n-v^n \\ e\geq 2}}p^e.\]
Assuming the $abc$-conjecture, for every $\epsilon>0$ we have
$B_n\ll_{\alpha,\epsilon} u^{\epsilon n}$.
\end{lemma}
\begin{proof}
Since
\[\rad(u^nv^n(u^n-v^n))\leq \rad(uv)A_nB_n^{1/2}=\rad(uv)\frac{u^n-v^n}{B_n^{1/2}}\leq\rad(uv)\frac{u^n}{B_n^{1/2}},\]
applying the $abc$-conjecture to
\[v^n+(u^n-v^n)=u^n\]
gives
\[u^n\ll_{\alpha,\delta}\left(\frac{u^n}{B_n^{1/2}}\right)^{1+\delta}\]
for every $\delta>0$. As a result, we have 
\[B_n^{(1+\delta)/2}\ll_{\alpha,\delta}u^{\delta n}\quad\implies\quad
B_n\ll_{\alpha,\delta}u^{\frac{2\delta}{1+\delta}n}.\]
The lemma follows by taking sufficiently small $\delta$. 
\end{proof}

\begin{proof}[Proof of Theorem \ref{logAm}(ii)]
Suppose the $abc$-conjecture holds. 
As before, we can assume that $\alpha=u/v$ with $u>v\geq 1$ and $\gcd(u,v)=1$. 
Suppose for contradiction that there exists $N$ such that $q_p(\alpha)\equiv 0\bmod p$ for each prime $p>N$ with $p\equiv 1\mod m$. 
\medskip

We consider primes $\ell>\max\{N,m,(u+v)^{\varphi(m)}\}$ and set 
\[W_\ell=\Phi_{m\ell}(u,v)=\frac{\Phi_m(u^\ell,v^\ell)}{\Phi_m(u,v)}.\] 
As in the previous part, every prime $p\mid W_\ell$ satisfies $p\nmid m\ell$ and $p\equiv 1\bmod m\ell$. In particular, 
\[p>\ell>N\qquad\text{and}\qquad p\equiv 1\mod m\]
so our assumption gives $q_p(\alpha)\equiv 0\bmod p$. Applying Lemma \ref{conglem} with $n=m\ell$, we obtain
\[\frac{W_\ell}{p}\equiv q_p(\alpha)v\Phi_{m\ell,Y}(u,v)\equiv 0\bmod p.\]
As a result, $p^2\mid W_\ell$ for every $p\mid W_\ell$ and $W_\ell$ is square-full (equivalently, powerful). 
Also, since $W_\ell\mid \left(u^{m\ell}-v^{m\ell}\right)$, we have 
\[W_\ell\mid B_{m\ell}\]
where $B_{m\ell}$ is the powerful part of $u^{m\ell}-v^{m\ell}$. Applying Lemma \ref{Silabc} with $\epsilon=\varphi(m)/2m$, we obtain
\[W_\ell\leq B_{m\ell}\ll_{\alpha,m}u^{\ell\varphi(m)/2}.\]
On the other hand, the cyclotomic estimates \eqref{trivest} give
\[W_\ell=\frac{\Phi_m(u^\ell,v^\ell)}{\Phi_m(u,v)}
\geq \frac{(u^\ell-v^\ell)^{\varphi(m)}}{(u+v)^{\varphi(m)}}
\geq \left(u^\ell\frac{1-v/u}{u+v}\right)^{\varphi(m)}
\gg_{\alpha,m} u^{\ell\varphi(m)}.\]
These upper and lower bounds for $W_\ell$ are incompatible for sufficiently large $\ell$ so we have a contradiction.
\end{proof}
\begin{remark}
The restriction $p\equiv 1\bmod m$ is intrinsic to the cyclotomic method presented here. 
Indeed, outside exceptional primes, a divisor $p\mid\Phi_{m\ell}(u,v)$ naturally lies in the class $p\equiv 1 \bmod m$. No analogous choice of these cyclotomic values appears to force a prescribed reduced class $p\equiv b\bmod m$ for $b\not\equiv 1\bmod m$ and the $abc$-conditional arguments of \cite{GRM13} and \cite{D19} have the same limitation. This reflects the familiar fact that the infinitude of primes $p\equiv 1\bmod m$ has an elementary cyclotomic proof, whereas Dirichlet's theorem on primes in arithmetic progressions lies much deeper.
\end{remark}


\subsection{Squares of finite logarithms in $\cA$} $\,$\\[3pt]
The condition $\log_\cA(\alpha)^2=c\in\mathbb{Q}^\times$ is equivalent to saying that for sufficiently large primes $p$, we have
\[q_p(\alpha)^2\equiv c\bmod p.\]
Since a non-square rational number is a quadratic non-residue modulo infinitely many primes, we must have $c=(a/b)^2$ for integers $a\neq 0$, $b\geq 1$, $\gcd(a,b)=1$. Then for every sufficiently large prime $p$, there is $\varepsilon_p\in\{\pm 1\}$ such that
\[q_p(\alpha)\equiv \varepsilon_p\frac{a}{b}\bmod p.\]
As before, for $\alpha=u/v\in\mathbb{Q}^\times\setminus\{\pm 1\}$, we need only consider $u>v\geq 1$ with $\gcd(u,v)=1$.
\smallskip

To obtain a contradiction, we will follow the same method as for Theorem \ref{logAm}. The essential differences are that $m$ will be chosen to be a suitable fixed auxiliary prime, the integer \eqref{Tell} will be replaced by a signed version involving $\varepsilon_p$, and $\ell$ will range over suitably large primes with $\ell\equiv 1\bmod m(m-1)$.


\begin{proof}[Proof of Theorem \ref{logsquare}] 
Suppose for contradiction there exists a bound $N>\max\{|a|,b\}$ such that for all primes $p>N$, we have
\[q_p(\alpha)\equiv\varepsilon_p\frac{a}{b}\bmod p \qquad\text{where $\varepsilon_p\in\{\pm 1\}$}.\] 
Fix an odd prime $m$ such that $m\nmid av(u-v)$ and set $w=\Phi_m(u,v)$.
For each prime $\ell>\max\{N,m,(u+v)^{m-1}\}$ with $\ell\equiv 1\bmod m(m-1)$, we set 
\[W_\ell=\Phi_{m\ell}(u,v)\qquad\text{and}\qquad R_\ell=\ell v^\ell\Phi_{m,Y}(u^\ell,v^\ell).\]
As before, any prime $p\mid W_\ell$ satisfies $p\equiv 1\mod m\ell$ and
\[w\frac{W_\ell}{p}\equiv q_p(\alpha)R_\ell\bmod p.\]
Furthermore, the right-hand side is non-zero modulo $p$ and $W_\ell$ is square-free. 
We define a signed version of \eqref{Tell} by 
\begin{equation*}
T_\ell^\varepsilon=bwS_\ell^\varepsilon-aR_\ell\in\mathbb{Z}\qquad\text{where}\qquad S_\ell^\varepsilon=\sum_{p\mid W_\ell}\varepsilon_p\frac{W_\ell}{p}\in\mathbb{Z}.
\end{equation*}
As in the proof of Theorem \ref{logAm}(i), we can conclude that $T_\ell^\varepsilon/W_\ell\in\mathbb{Z}$. 
Indeed, for $p\mid W_\ell$, we have $S_\ell^\varepsilon\equiv \varepsilon_pW_\ell/p\bmod p$ and so 
\[T_\ell^\varepsilon\equiv b\varepsilon_pq_p(\alpha)R_\ell-aR_\ell\equiv 0\bmod p.\]
Similarly, we can conclude that $T_\ell^\varepsilon/W_\ell\to 0$ as $\ell\to\infty$. 
Indeed, this follows since $\left|R_\ell/W_\ell\right|\to 0$ as before, and Lemma \ref{limlem} gives
\[\left|\frac{S_\ell^\varepsilon}{W_\ell}\right|\leq \sum_{p\mid W_\ell}\frac{1}{p}\longrightarrow 0.\]
The only remaining step is to show that $T_\ell^\varepsilon/W_\ell$ is non-zero. 
Set 
\[E_\ell=\sum_{p\mid W_\ell}\varepsilon_p.\]
As before, we have $S_\ell^\varepsilon\equiv E_\ell\bmod \ell$, $R_\ell\equiv 0\bmod \ell$ and hence 
\[T_\ell^\varepsilon\equiv bwE_\ell\bmod \ell.\]
We have chosen $\ell$ so that $1\leq b<\ell$ and $1\leq w<\ell$ so $bw\not\equiv 0\bmod\ell$. 
We also chose $\ell>(u+v)^{m-1}$ and hence by Lemma \ref{limlem},
\[\left|E_\ell\right|\leq r_{m\ell}<\frac{(\ell-1)(m-1)\log(u+v)}{\log(m\ell)}<\ell-1.\]
Thus, if $E_\ell\neq 0$, we deduce that $T_\ell^\varepsilon\not\equiv 0\bmod \ell$.
\medskip

On the other hand, suppose $E_\ell=0$. 
Since every prime $p\mid W_\ell$ satisfies $p\equiv 1 \bmod m$ and $W_\ell$ is square-free, we have $W_\ell/p\equiv 1\bmod m$. Hence $S_\ell^\varepsilon\equiv E_\ell\equiv 0\bmod m$ and
\[T_\ell^\varepsilon=bwS_\ell^\varepsilon-aR_\ell\equiv -aR_\ell\bmod m.\]
Now $\Phi_m(X,Y)\equiv (X-Y)^{m-1}\bmod m$ and so $\Phi_{m,Y}(X,Y)\equiv (X-Y)^{m-2}\bmod m$. 
Furthermore, we are letting $\ell$ range over the infinitely many primes of the form
\[\ell\equiv 1\bmod m(m-1).\] 
Thus $\ell\equiv 1\bmod m$ and $\ell\equiv 1\bmod m-1$ so $u^\ell\equiv u\bmod m$ and $v^\ell\equiv v\bmod m$. 
In particular, we now have
\[R_\ell=\ell v^\ell\Phi_{m,Y}(u^\ell,v^\ell)\equiv v(u-v)^{m-2}\not\equiv 0\bmod m.\]
Thus, if $E_\ell=0$, we deduce that $T_\ell^\varepsilon\not\equiv 0\bmod m$.
\medskip

In either case, we find $T_\ell^\varepsilon/W_\ell$ is a non-zero integer which tends to zero as $\ell\to\infty$, giving the required contradiction. 
This completes the proof of Theorem \ref{logsquare}(i). 
Finally, part (ii) follows immediately from Theorem \ref{MSS}(ii), since $\cA$ is a reduced ring: for any ${\bf t}\in\cA$, if ${\bf t}^2=0$ then ${\bf t}=0$.
\end{proof}


\begin{remark}
There is an alternative finite logarithm that can be considered in $\cA$, as well as finite analogues of higher polylogarithms. 
For prime $p$ and integer $d\geq 1$, the {\it finite $d$-logarithm} is
\[\LL_{d}(X)=\sum_{k=1}^{p-1}\frac{X^k}{k^d}.\]
For $d=1,2$, these are related to Fermat quotients modulo $p$ and in particular (see e.g. \cite{Gr04}), for $p>3$ we have simple identities
\[\LL_{1}(2)\equiv -2q_p(2)\bmod p\qquad\text{and}\qquad\LL_{2}(2)\equiv -q_p(2)^2\bmod p.\]
In view of Theorems \ref{MSS} and \ref{logsquare}, we deduce $\left(\LL_{1}(2)\bmod p\right)_p$ and $\left(\LL_{2}(2)\bmod p\right)_p$ are not in $\mathbb{Q}^\times\subset\cA$ unconditionally and, assuming the $abc$-conjecture, are not rational elements of $\cA$.
\end{remark}

\begin{thebibliography}{99}
\bibitem[D19]{D19} \textsc{Y. Ding}, Non-{W}ieferich primes under the $abc$ conjecture, \emph{C. R. Acad. Sci. Paris, Ser. I} \textbf{357} (2019), 483--486.
\bibitem[GRM13]{GRM13} \textsc{H. Graves} and \textsc{M. Ram Murty}, The abc conjecture and non-{W}ieferich primes in arithmetic progressions, \emph{J. Number Theory} \textbf{133} (2013), no.~6, 1809--1813.
\bibitem[Gr04]{Gr04} \textsc{A. Granville}, The square of the Fermat quotient, \emph{Integers} \textbf{4} (2004), \#A22.
\bibitem[KZ]{KZ} \textsc{M. Kaneko} and \textsc{D. Zagier}, Finite multiple zeta values, \emph{to appear in the Proceedings of the 17th MSJ-SI conference on Modular forms and Multiple Zeta values.}
\bibitem[LZ25]{LZ25} \textsc{F. Luca} and \textsc{W. Zudilin}, Irrationality and transcendence questions in the `poor man's ad\`ele ring', \emph{Ramanujan J.} \textbf{67} (2025), no.~4, Paper No. 88.
\bibitem[MS26]{MS26} \textsc{T. Matsusaka} and \textsc{S. Seki}, Some results on naive transcendence in the ring of integers modulo infinitely large primes, \emph{Preprint} \href{https://arxiv.org/abs/2604.25566}{\texttt{arXiv:2604.25566 [math.NT]}} (2026).
\bibitem[RTTY24]{RTTY24} \textsc{J. Rosen}, \textsc{Y. Takeyama}, \textsc{K. Tasaka} and \textsc{S. Yamamoto}, The ring of finite algebraic numbers and its application to the law of decomposition of primes, \emph{J. Number Theory} \textbf{263} (2024), 335--365.
\bibitem[Ro20]{Ro20} \textsc{J. Rosen}, A finite analogue of the ring of algebraic numbers, \emph{J. Number Theory} \textbf{208} (2020), 59--71.
\bibitem[Si88]{Si88} \textsc{J.\,H. Silverman}, Wieferich's criterion and the $abc$-conjecture, \emph{J. Number Theory} \textbf{30} (1988), no.~2, 226–237.
\end{thebibliography}
\end{document}
